% Source: https://artofproblemsolving.com/community/c3h1523697_problem_solving_finding_minimum_and_maximum
% Stable identity: aops-topic-1523697-q01
Let $g$ and $o$ be the litres of grape juice and orange juice used, and $T=g+o$ the total volume. The pure juice is
\[
0.2g+0.5o=24 .
\]
The juice content $\dfrac{24}{T}$ lies between $0.25$ and $0.375$ exactly when $\dfrac{24}{0.375}\le T\le\dfrac{24}{0.25}$, that is $64\le T\le96$.

Substituting $o=T-g$ gives $0.5T-0.3g=24$, so
\[
g=\dfrac{0.5T-24}{0.3}=\dfrac{5T-240}{3},
\]
which increases with $T$. At $T=64$: $g=\dfrac{80}{3}=26\dfrac{2}{3}$ litres (and $o=\dfrac{112}{3}$). At $T=96$: $g=80$ litres (and $o=16$). Both amounts are nonnegative throughout.

So the least amount is $\boxed{26\dfrac{2}{3}\text{ litres}}$ and the greatest is $\boxed{80\text{ litres}}$.
